\hypertarget{chapter-18-c14-templates-and-metaprogramming}{%
\section{CHAPTER 18: C14 TEMPLATES AND
METAPROGRAMMING}\label{chapter-18-c14-templates-and-metaprogramming}}

\hypertarget{c14-templates-metaprogramming}{%
\section{C++14 TEMPLATES \&
METAPROGRAMMING}\label{c14-templates-metaprogramming}}

C++14 simplified template metaprogramming (TMP) by introducing variable
templates and utility classes that replaced complex, recursive
boilerplate with cleaner, more intuitive syntax.

\hypertarget{variable-templates}{%
\subsection{1. Variable Templates}\label{variable-templates}}

Before C++14, if you wanted a templated constant (like
\passthrough{\lstinline!pi!}), you had to wrap it in a
\passthrough{\lstinline!struct!} or a
\passthrough{\lstinline!constexpr!} function. Variable templates allow
direct templating of variables.

\hypertarget{mathematical-constants-and-type-traits}{%
\subsubsection{1.1 Mathematical Constants and Type
Traits}\label{mathematical-constants-and-type-traits}}

This feature is heavily used in the standard library and mathematical
libraries.

\begin{lstlisting}[language={C++}]
template<typename T>
constexpr T pi = T(3.1415926535897932385L);

// Usage
float  f_pi = pi<float>;
double d_pi = pi<double>;

// Type traits (Internal simplification)
template <typename T>
constexpr bool is_floating_point_v = std::is_floating_point<T>::value;
\end{lstlisting}

\hypertarget{professional-note-_v-suffixes}{%
\subsubsection{\texorpdfstring{1.2 Professional Note: \texttt{\_v}
Suffixes}{1.2 Professional Note: \_v Suffixes}}\label{professional-note-_v-suffixes}}

C++14 (and later C++17) introduced \passthrough{\lstinline!\_v!} aliases
for most type traits. Instead of writing
\passthrough{\lstinline!std::is\_integral<T>::value!}, you can write
\passthrough{\lstinline!std::is\_integral\_v<T>!}. This reduces noise in
complex template expressions.

\hypertarget{stdinteger_sequence-the-indices-trick}{%
\subsection{\texorpdfstring{2. \texttt{std::integer\_sequence} \& The
Indices
Trick}{2. std::integer\_sequence \& The Indices Trick}}\label{stdinteger_sequence-the-indices-trick}}

Meta-programming often involves working with variadic templates and
tuples. \passthrough{\lstinline!std::integer\_sequence!} provides a way
to generate a sequence of integers at compile-time.

\hypertarget{unpacking-a-tuple}{%
\subsubsection{2.1 Unpacking a Tuple}\label{unpacking-a-tuple}}

The ``Indices Trick'' is the classic use case: converting a
\passthrough{\lstinline!std::tuple!} into a pack of arguments for a
function.

\begin{lstlisting}[language={C++}]
template<typename F, typename Tuple, std::size_t... I>
auto apply_impl(F f, Tuple&& t, std::index_sequence<I...>) {
    return f(std::get<I>(std::forward<Tuple>(t))...);
}

template<typename F, typename Tuple>
auto apply(F f, Tuple&& t) {
    using Indices = std::make_index_sequence<std::tuple_size_v<std::decay_t<Tuple>>>;
    return apply_impl(f, std::forward<Tuple>(t), Indices{});
}

// Result: apply(func, make_tuple(1, 2)) calls func(1, 2)
\end{lstlisting}

\textbf{Godhood Insight:} \passthrough{\lstinline!std::index\_sequence!}
(an alias for
\passthrough{\lstinline!std::integer\_sequence<size\_t, ...>!}) is the
glue that connects the ``Value World'' (Tuples) to the ``Pack World''
(Variadic Templates).

\hypertarget{alias-templates-and-_t-suffixes}{%
\subsection{\texorpdfstring{3. Alias Templates and \texttt{\_t}
Suffixes}{3. Alias Templates and \_t Suffixes}}\label{alias-templates-and-_t-suffixes}}

C++14 introduced alias templates for all type traits in
\passthrough{\lstinline!<type\_traits>!}.

\hypertarget{reducing-typename-...type-boilerplate}{%
\subsubsection{\texorpdfstring{3.1 Reducing \texttt{typename\ ...::type}
Boilerplate}{3.1 Reducing typename ...::type Boilerplate}}\label{reducing-typename-...type-boilerplate}}

In C++11, using a trait that returns a type required the
\passthrough{\lstinline!typename!} keyword and the
\passthrough{\lstinline!::type!} suffix. C++14 added
\passthrough{\lstinline!\_t!} versions.

\begin{lstlisting}[language={C++}]
// C++11 (Verbosely painful)
typename std::enable_if<Condition, T>::type

// C++14 (Clean and readable)
std::enable_if_t<Condition, T>

// Example: SFINAE made easy
template<typename T, typename = std::enable_if_t<std::is_integral_v<T>>>
void only_integers(T val) { /* ... */ }
\end{lstlisting}

\textbf{Professional Note:} Always prefer the
\passthrough{\lstinline!\_t!} and \passthrough{\lstinline!\_v!} versions
in modern C++. They are not just shorter; they are conceptually cleaner
because they treat the trait as a function that returns a type/value
directly.
